\section{The original theta integral: real derivatives, uniform growth, and lowering}
\label{sec:theta-extension-v034}

\status{TAE1--TAE13 independently reconstruct the real analytic and
lowering block of the original base theta extension.  They apply to the
same full Weil datum, theta functions, frames, coefficient space, Haar
measure and fundamental domain already used in BTH, FIA and CTC.
The remaining common polynomial identification, lowest-type
multiplicity-one, boundedness and Hilbert realization are stated
separately at the end.  No new Lean replay or category-A application is
claimed.}

Fix the original number-field datum \(D\), standard character
\(\psi\) with \(\mathrm{IsStdChar}(D,\psi,-1)\), Hermitian spaces
\(S,T\), infinity chart \(\theta\), and full action
\[
 \omega:G_S(\mathbb A_F)\times G_T(\mathbb A_F)
       \longrightarrow\operatorname{End}_{\mathbb C}
                    \bigl(\mathbb C^{W}\bigr),
 \qquad W=\mathrm{WO}(S,T).
 \tag{TAE.1}
\]
The premise on this action is the literal
\(\mathrm{IsWeilDatumO}(\psi,S,T,\chi_S,\chi_T,\omega)\).
Write \(\Theta\) for its original space
\(\mathrm{IsThetaFunO}(\psi,S,T)\).  The real Hom space is
\[
 V_{\mathbb R}=\mathrm{RW}(T,S)
       =\prod_{\tau}\operatorname{Mat}_{r_T,r_S}(\mathbb C),
 \qquad N=\dim_{\mathbb R}V_{\mathbb R}.
 \tag{TAE.2}
\]
The inverse infinity chart gives mutually compatible real inclusion
\(\iota:V_{\mathbb R}\to W\) and projection \(\pi:W\to V_{\mathbb R}\).
These are the actual \(\mathrm{iotaW}\) and \(\mathrm{piW}\).
The real Heisenberg action is
\(\rho(v)=\mathrm{heisO}(\psi,S,T,\iota(v))\).

CW12 identifies every \(f\in\Theta\) with a smooth vector in the Banach
space \(\mathcal B\) of bounded continuous functions satisfying the
original rational quasi-periodicity, equipped with the supremum norm.
We use its smooth topology
\[
 p_{\mathbf v}(f)=
  \|D_{v_1}\cdots D_{v_k}f\|_{\infty},
 \qquad
 D_vf=\left.\frac{d}{dt}\rho(tv)f\right|_{t=0}.
 \tag{TAE.3}
\]
The empty word is included.  Every \(D_v\) preserves \(\Theta\), and
all real derivatives preserve any fixed finite Heisenberg level.
These facts concern the genuine smooth domain; they do not assert
continuity of the ambient endomorphism on all functions.

The theta kernel and the source integral will always mean
\[
 K_\omega(h,g;f)=(\omega(h,g)f)(0),\qquad
 \Phi_L(g)=\int_{\mathcal F_S}
        K_\omega(h,g;L(F(h)))\,d\nu(h).
 \tag{TAE.4}
\]
Here \(E\) is a finite-dimensional complex normed space,
\(L:E\to\Theta\) is complex linear, and
\(F:G_S(\mathbb A_F)\to E\) is continuous and left
\(G_S(F)\)-invariant.  Whenever the integral is used, its original
source quotient has a compact set of representatives \(C_S\).
For the supplied base bundle this is exactly the finite class
presentation established in BTH3, together with compactness of the
definite real group and of the finite level.  Thus no new general
class-number assertion is being imported.

\subsection*{The real action on its genuine smooth domain}

% BEGIN TAE1
\paragraph{TAE1: the explicit Gaussian oscillator action is smooth in all real parameters.}
\label{claim:TAE1}
Use CW8--CW11's real coordinates: \(V\) is Euclidean,
\(J^2=-1\) is orthogonal, and
\[
 \rho(u)\rho(v)=e^{\pi i\beta(u,v)}\rho(u+v),
 \qquad \beta(u,v)=\langle u,Jv\rangle .
 \tag{TAE.5}
\]
Let \(g\) vary in the unitary matrix group preserving the same
Hermitian form.  CW10's operator
\(\Omega(g)\) acts smoothly on each smooth vector, in every seminorm
(TAE.3).  Its derivatives preserve the smooth-vector space.  On a
compact parameter set each derivative is bounded by a finite sum of
seminorms of the input.

\emph{Proof.}
Put \(G(u)=e^{-\pi\|u\|^2/2}\),
\[
 P=\int_VG(u)\rho(u)\,du,\qquad
 a_f(v)=P\rho(-v)f .
 \tag{TAE.6}
\]
The Gaussian and the Euclidean Lebesgue measure are precisely those
of CW8.  That proof gives
\[
 a_f(v)=\rho(-v)\widehat{(G(\cdot)\rho(\cdot)f)}(Jv),
 \qquad
 \int_V(1+\|v\|)^q\|a_f(v)\|\,dv<\infty
 \quad(q\ge0).
 \tag{TAE.7}
\]
The integral in (TAE.7) is controlled by finitely many smooth
seminorms of \(f\).  Indeed differentiate
\(G(u)\rho(u)f\) in a fixed real basis.  Each derivative is a finite
sum of a polynomial times \(G(u)\rho(u)D_{\mathbf e}f\), by the
Heisenberg law.  Its \(L^1\) norm is therefore bounded by a finite
sum of the \(p_{\mathbf e}(f)\).  Fourier integration by parts,
using more than \(q+\dim V\) coordinate derivatives, bounds its
Fourier transform by a constant times
\((1+\|v\|)^{-q-\dim V-1}\).  This proves the claimed quantitative
version of (TAE.7).

We also need a uniform estimate on derivatives of \(P\), not just
pointwise differentiability.  Changing variables in
\[
 \rho(tz)Pf
  =\int_VG(u-tz)e^{\pi i\beta(tz,u)}\rho(u)f\,du
 \tag{TAE.8}
\]
shows that \(D_zP\) is a bounded operator on \(\mathcal B\):
its kernel is a degree-one polynomial times a Gaussian.
Repeating the argument proves that
\(D_{z_1}\cdots D_{z_k}P\) is bounded, with norm at most
\(C_k\prod_j\|z_j\|\).  All differentiations of (TAE.8) are
dominated by an integrable polynomial times a Gaussian, uniformly
for \(t\) in a compact interval.  Since \(P^2=P\), these estimates
apply to \(a_f(v)\), with its own norm on the right.

Expanding the two projectors in CW10 and using covariance gives
the following particularly useful version of its defining integral:
\[
 \Omega(g)f=c(g)\int_V\int_V
   G(u)e^{\pi i\beta(v,u)}
           \rho(g(v+u))a_f(v)\,du\,dv .
 \tag{TAE.9}
\]
Here \(c(g)=\det(p g+q)\) in the signature coordinates of CW9;
after a fixed Hermitian congruence it is the same determinant
polynomial of the conjugated matrix.  The phase in (TAE.9) is
independent of \(g\), since \(g\) preserves \(\beta\).

For a smooth curve \(y(t)\) and a smooth vector \(a\), the
Heisenberg law yields
\[
 \frac{d}{dt}\rho(y(t))a
 =\rho(y(t))
       \bigl(D_{y'(t)}a-\pi i\beta(y(t),y'(t))a\bigr).
 \tag{TAE.10}
\]
Apply this with \(y(t)=g(t)(v+u)\) and \(a=a_f(v)\).
All derivatives of \(g\) are bounded on a fixed compact parameter
set.  Iterating (TAE.10), and then applying any fixed word of
left Heisenberg derivatives, bounds the resulting integrand by
\[
 C(1+\|u\|+\|v\|)^q\,G(u)\,\|a_f(v)\|.
 \tag{TAE.11}
\]
The derivatives of \(a_f(v)\) which occur here are bounded using
(TAE.8), because \(a_f(v)=Pa_f(v)\).
The function in (TAE.11) is integrable by (TAE.7) and the Gaussian
moments in \(u\).  Dominated differentiation in the Banach norm
therefore applies at every order.  It also proves continuity of
each derivative in every seminorm (TAE.3).

For clarity, this argument is valid even though a Heisenberg
derivative is an unbounded operator on \(\mathcal B\).
One can apply a finite Heisenberg difference quotient to (TAE.9)
first, dominate it by (TAE.11), and then pass to the limit.  The
limit is the derivative asserted above.  Equivalently the
generator of an isometric one-parameter group is closed, so the
same dominated estimates identify its value on the integral.
Thus no exchange of an unbounded operator and an integral is
assumed.

In the original adelic model, a fixed finite Heisenberg
translation commutes with every real \(\rho(u)\).  It therefore
commutes with \(P\), with (TAE.9), and with all the derivatives
just obtained.  A finite open level fixing \(f\) fixes all those
derivatives.  CW12's converse now makes them original theta
functions.  This proves the assertion on exactly the space that
will be used below.
% END TAE1

% BEGIN TAE2
\paragraph{TAE2: the continuous scalar twists have genuine real derivatives.}
\label{claim:TAE2}
A continuous character \(s:\mathbb R\to\mathbb C^\times\) with
\(s(t+u)=s(t)s(u)\) is smooth and has the form
\(s(t)=e^{a t}\).  A continuous character of a finite-dimensional
matrix Lie group is smooth on that group.  In particular the
determinant-root characters in CW24 are smooth on the real
factor of the original adelic group.

\emph{Proof.}
Let \(A(t)=\int_0^t s(u)\,du\).  Then \(A'=s\), and
\(A(\epsilon)\ne0\) for some \(\epsilon\), since
\(A'(0)=s(0)=1\).  Translation of this ordinary integral gives
\[
 s(t)=\frac{A(t+\epsilon)-A(t)}{A(\epsilon)}.
 \tag{TAE.12}
\]
The right side is continuously differentiable.  Differentiating
the character law at its second variable zero gives
\(s'(t)=s(t)s'(0)\).  Hence
\(e^{-s'(0)t}s(t)\) has zero derivative and value one at zero.
This proves the formula and all higher differentiability.
If the character is unitary, \(a\) is purely imaginary, but this
extra fact is not needed for differentiability.

Choose a real basis \(X_1,\ldots,X_d\) of a matrix Lie algebra.
The product
\(\exp(t_1X_1)\cdots\exp(t_dX_d)\) is a smooth local chart at the
identity: its differential there is the basis isomorphism, so the
inverse function theorem applies.  A character on this chart
is the product of the smooth one-variable characters just proved.
Translation gives smoothness near every point.  CW24 proves that
its determinant-root twists are continuous characters on the
full original adelic groups.  Restriction to the real factors
therefore satisfies precisely these hypotheses.
% END TAE2

% BEGIN TAE3
\paragraph{TAE3: the original action has the required generator, commutators, and translated point derivatives.}
\label{claim:TAE3}
For every real place \(\tau\) and
\[
 X^*H_{T,\tau}+H_{T,\tau}X=0
 \tag{TAE.13}
\]
there is a complex-linear map \(\mathcal G_{\tau,X}:\Theta\to\Theta\)
such that, for every \(f\in\Theta\), every adelic pair \(p\), and
every \(w\in W\),
\[
 \left.\frac{d}{dt}
  \bigl(\omega(p(1,\mathrm{oneP}_\tau(X,t)))f\bigr)(w)
       \right|_{t=0}
       =(\omega(p)\mathcal G_{\tau,X}f)(w).
 \tag{TAE.14}
\]
Writing \(D_{\tau',u}\) for the original \(\mathrm{dRho}\), one has
\[
 [\mathcal G_{\tau,X},D_{\tau,u}]f=D_{\tau,Xu}f,
 \qquad
 [\mathcal G_{\tau,X},D_{\tau',u}]f=0\quad(\tau'\ne\tau).
 \tag{TAE.15}
\]
Extending \(\mathcal G_{\tau,X}\) linearly off \(\Theta\) gives
the literal ambient endomorphism in the original
\texttt{exists\_\allowbreak generator} statement.

\emph{Proof.}
CW24 constructs a full datum with the same
\(D,\psi,S,T,\chi_S,\chi_T\).  BTH14 identifies it with the
specified \(\omega\) on every theta function.  On the real
factor this constructed action is
\[
                  \omega(\gamma)f=\zeta(\gamma)\Omega(m(\gamma))f,
 \tag{TAE.16}
\]
where \(\zeta\) is the continuous unitary determinant-root
character of CW24 and
\[
 m(g,h)=\operatorname{diag}_{\tau'}
       \bigl(h_{\tau'}\otimes(g_{\tau'}^{-1})^{\mathsf t}\bigr)
 \tag{TAE.17}
\]
is the original action on the real Hom space.  Neither a new
character nor a different action is chosen in (TAE.16).
The finite factor of a real-supported element is the identity.
TAE1 and TAE2 therefore make (TAE.16) smooth in the topology
(TAE.3).  Its derivative along \((1,\mathrm{oneP}_\tau(X,t))\)
defines \(\mathcal G_{\tau,X}\).  Linearity in \(f\) follows from
linearity of the action.  Preservation of the real smooth
domain and of a finite level follows from TAE1; scalar
differentiation does not change a finite level.  Thus its
values are in \(\Theta\).  It is also real linear in \(X\),
being the differential of the smooth real group action.

The group covariance law gives, on \(\Theta\),
\[
 \omega(1,\mathrm{oneP}_\tau(X,t))D_{\tau',u}
 =D_{\tau',\,e^{tX}u}
       \omega(1,\mathrm{oneP}_\tau(X,t))
       \quad(\tau'=\tau);
 \tag{TAE.18}
\]
if \(\tau'\ne\tau\), its right-hand direction is just \(u\).
To justify (TAE.18), first differentiate the Heisenberg
covariance law in the real shift parameter.  Both sides are
smooth in the topology (TAE.3) by TAE1, so the differentiation
is legitimate.  Differentiating once more in \(t\) at zero
gives exactly (TAE.15), including its plus sign.

There is a separate point concerning the arbitrary adelic
left factor \(p\) in (TAE.14).  We do not assume that
\(\omega(p)\) is bounded in the supremum norm.
Set
\[
 u(t)=\omega(\gamma(t))f,\quad
 A=\mathcal G_{\tau,X}f,\quad
 v(t)=\omega(p\gamma(t))f,
 \qquad \gamma(t)=(1,\mathrm{oneP}_\tau(X,t)).
 \tag{TAE.19}
\]
The first curve is differentiable in \(\mathcal B\).
The real subgroup is normal in the adelic group, and
\(p\gamma(t)p^{-1}\) is a smooth real matrix curve.
Applying (TAE.16) to this curve and to \(\omega(p)f\in\Theta\)
shows that \(v(t)\) is also differentiable in \(\mathcal B\);
write its derivative \(B\).

Choose an additive Haar measure on \(W\), normalized so that
the compact rational quotient has volume one, and a rational
fundamental domain \(\mathcal F_W\).  The map
\(\mathcal B\to L^2(\mathcal F_W)\) is bounded.  The original
\(\mathrm{om\_unit}\) gives
\[
 \|v(t)-v(0)-t\,\omega(p)A\|_{L^2}
       =\|u(t)-u(0)-tA\|_{L^2}=o(|t|).
 \tag{TAE.20}
\]
All vectors inside this equality are theta functions, so that
the stated isometry applies.  Uniqueness of an \(L^2\)
derivative gives \(B=\omega(p)A\) in \(L^2\).
Their difference is continuous and its norm is rationally
periodic.  A nonzero value would give a positive-measure open
set of positive norm in the quotient; hence equality in
\(L^2\) implies equality at every point of \(W\).
Point evaluation on \(\mathcal B\) is bounded, so (TAE.14)
follows.

Finally choose a complex vector-space complement to \(\Theta\)
in the space of all functions and let the generator be zero
on that complement.  This gives an ambient complex-linear
endomorphism.  All assertions involve theta functions and
their Heisenberg derivatives, so the extension changes none
of them.  This proves exactly the original generator
contract, without leaving a real oscillator differentiability
theorem as an additional input.
% END TAE3

\subsection*{One growth exponent for every derivative word}

% BEGIN TAE4
\paragraph{TAE4: the finite-level real-cube bound is uniform along the real group.}
\label{claim:TAE4}
Fix a basis \(e_1,\ldots,e_N\) of \(V_{\mathbb R}\), an open
finite Heisenberg level \(U\), a target point \(x\), and a compact
source set \(C\).  There is \(b_0\ge0\), independent of
\(h\in C\), \(g\in G_{T,\infty}\), and \(\varphi\in\Theta\)
fixed at level \(U\), such that
\[
 \int_{[0,1]^N}
   \left|\,\omega(h,xg)\varphi
          \left(\iota\left(\sum_j t_je_j\right)\right)\right|^2dt
             \le b_0\|\varphi\|_\infty^2.
 \tag{TAE.21}
\]
The same constant applies to every real Heisenberg derivative
of every such \(\varphi\).

\emph{Proof.}
The action of \((h,xg)^{-1}\) on a purely finite shift depends
only on \(h,x\); the real-supported \(g\) has trivial finite
component.  Continuity of this linear action and compactness
of \(C\) provide a common compact open neighborhood
\(U_0\) in the finite Hom space whose image lies in \(U\) for
every \(h\in C\).  Here one first chooses a neighborhood for
each \(h\), covers \(C\) by finitely many corresponding
neighborhoods, and intersects the finitely many finite
neighborhoods.  The full covariance law consequently fixes
\(\omega(h,xg)\varphi\) under the Heisenberg shifts in \(U_0\).
Their phases have modulus one.  Thus the squared norm at
\(\iota(v)+u\), \(u\in U_0\), equals its squared norm at
\(\iota(v)\).

The continuous splitting \(W=V_{\mathbb R}\oplus W_f\)
identifies Haar measure with a positive scalar multiple of
the product Haar measures.  The compact set
\[
 K_0=\left\{\iota\left(\sum_jt_je_j\right)+u:
                         t\in[0,1]^N,\ u\in U_0\right\}
 \tag{TAE.22}
\]
therefore converts the left side of (TAE.21) to a fixed
positive multiple of the integral of
\(|\omega(h,xg)\varphi|^2\) over \(K_0\).
This is the explicit real-slice measure comparison; restricting
an adelic integral to a measure-zero real slice would not give
it without the common finite level.

Let \(\Gamma=W(F)\), a discrete cocompact additive lattice.
Each coset meets \(K_0\) with multiplicity at most
\[
 M_0=\#\bigl(\Gamma\cap(K_0-K_0)\bigr)<\infty.
 \tag{TAE.23}
\]
Indeed, relative to one such intersection point, every other
rational translate lies in the displayed finite set.
Unfolding a nonnegative rationally periodic function over a
fundamental domain consequently bounds its \(K_0\)-integral
by \(M_0\) times its quotient integral.
The norm of a theta function is rationally periodic, and
\(\mathrm{om\_unit}\) identifies the quotient integral of
\(|\omega(h,xg)\varphi|^2\) with that of \(|\varphi|^2\).
Normalize the latter quotient measure to have volume one;
it is then at most \(\|\varphi\|_\infty^2\).
Combining these bounds gives (TAE.21).

Real and finite Heisenberg shifts commute.  Differentiating
that equality proves that every real derivative of
\(\varphi\) is fixed at the same \(U\).  Nothing in the
construction of \(U_0,K_0,M_0,b_0\) depends on the particular
vector, proving the final assertion.
% END TAE4

% BEGIN TAE5
\paragraph{TAE5: the original theta kernel has exponent \(N=\dim_{\mathbb R}\mathrm{RW}(T,S)\).}
\label{claim:TAE5}
For every \(f\in\Theta\), target point \(x\), and compact
source set \(C\), there is \(c_f\ge0\) such that
\[
 |K_\omega(h,xg;f)|
       \le c_f\,H_T(g)^N
       \quad(h\in C,\ g\in G_{T,\infty}),
 \qquad H_T(g)=\mathrm{normInf}(\theta,T,g).
 \tag{TAE.24}
\]
The exponent is independent of \(f,x,C\).

\emph{Proof.}
Use the supremum norm on the coordinates of \(V_{\mathbb R}\).
The original height is exactly
\[
 H_T(g)=\sum_{\tau,a,b}
       \left(|(g_\tau)_{ab}|+|(g_\tau^{-1})_{ab}|\right).
 \tag{TAE.25}
\]
For the original real Hom action write
\[
 M(h,k)u=(k_\tau u_\tau h_\tau^{-1})_\tau .
 \tag{TAE.26}
\]
The full Heisenberg covariance, differentiated as in TAE3,
gives
\[
 D_u\,\omega(h,k)f
        =\omega(h,k)\,D_{M(h^{-1},k^{-1})u}f .
 \tag{TAE.27}
\]
The same \(L^2\) difference-quotient argument as (TAE.20)
justifies this formula for an arbitrary adelic pair; all
Heisenberg derivatives exist in \(\mathcal B\) by CW12.

Matrix multiplication and the triangle inequality give the
precise height comparison
\[
 \|M(h,(xg)^{-1})u\|
       \le \|M(h,x^{-1})u\|\,H_T(g).
 \tag{TAE.28}
\]
To verify it, the left matrix is
\(g_\tau^{-1}(x_\tau^{-1}u_\tau h_\tau^{-1})\).
Each row sum of \(|g_\tau^{-1}|\) is bounded by (TAE.25),
and each entry of the other factor is bounded by its
supremum norm.  No changed height or hidden exponent is used.
Compactness of \(C\) gives a single \(\kappa\ge0\) with
\[
 \|M(h^{-1},x^{-1})e_j\|\le\kappa
                       \quad(h\in C,\ 1\le j\le N).
 \tag{TAE.29}
\]

We record an elementary word estimate.  For each fixed
\(f\in\Theta\) there is \(A_f\ge0\) such that
\[
 \|D_{u_k}\cdots D_{u_1}f\|_\infty
        \le A_f\rho^N
 \quad\left(k\le N,\ \rho\ge1,\ \|u_j\|\le\rho\right).
 \tag{TAE.30}
\]
Expand every \(u_j\) in the fixed real basis.  The coordinate
functionals have bounded norm, and \(D_u\) is real linear in
\(u\).  The expansion is a finite sum of the finitely many
basis words of lengths at most \(N\), with coefficient bounds
of the form \(C^k\rho^k\).  Each basis word is a bounded theta
function by CW12.  Summing their bounds and using
\(\rho^k\le\rho^N\) proves (TAE.30), including the empty word.

For a real direction \(e_j\), the one-dimensional fundamental
theorem of calculus and Cauchy--Schwarz imply
\[
 |q(0)|^2\le
       2\int_0^1|q(t)|^2\,dt+
       2\int_0^1|q'(t)|^2\,dt .
 \tag{TAE.31}
\]
Indeed \(q(0)=q(t)-\int_0^tq'(r)\,dr\); square, integrate
over \(t\), and enlarge the derivative integral to \([0,1]\).
Apply this to a Heisenberg orbit evaluated at a point.
The unit-modulus phase shows that its two terms are the
ordinary squared norms of the vector and its covariant
derivative on the translated real segment.  Iterating along
the \(N\) coordinate directions gives \(4^N\) times a common
upper bound for all real-cube integrals of derivative words
of length at most \(N\).  This is also an immediate induction:
at each step there are two terms, each with coefficient two.

Apply this inequality to \(\omega(h,xg)f\) at zero.
By (TAE.27)--(TAE.29), every derivative that appears is
\(\omega(h,xg)\varphi\), where \(\varphi\) is a word on \(f\)
of length at most \(N\) whose directions have norms at most
\(\kappa H_T(g)\).  All such words have the same finite
level.  TAE4 and (TAE.30), with
\(\rho=\max(1,\kappa H_T(g))\), therefore give
\[
 |K_\omega(h,xg;f)|^2
        \le 4^N b_0 A_f^2\rho^{2N}.
 \tag{TAE.32}
\]
If \(N>0\), then \(H_T(g)\ge1\).  To see this, select one
nonempty real matrix block and an index \(a\); the identity
\(g_\tau g_\tau^{-1}=1\) gives
\(1\le\sum_b |(g_\tau)_{ab}||(g_\tau^{-1})_{ba}|
 \le H_T(g)^2\).
Hence
\(\rho^N\le\max(1,\kappa)^N H_T(g)^N\).
Taking square roots in (TAE.32) proves (TAE.24), with
\[
       c_f=2^N\sqrt{b_0}\,A_f\max(1,\kappa)^N .
 \tag{TAE.33}
\]
If \(N=0\), the cube has one point, all words are empty, and
\(H_T(g)^0=1\); (TAE.21) proves the same assertion directly.
Thus the degenerate case needs no positive-height assumption.
% END TAE5

\subsection*{Differentiating the original definite-source integral}

% BEGIN TAE6
\paragraph{TAE6: compact representatives give dominated differentiation of the original integral.}
\label{claim:TAE6}
For the data of (TAE.4), \(\Phi_L\) is continuous and left
\(G_T(F)\)-invariant.  For every \(\tau,X\) satisfying
(TAE.13), put \(L'=\mathcal G_{\tau,X}\circ L\).  Then
\[
 R_{\tau,X}\Phi_L=\Phi_{L'},\qquad
 \left.\frac{d}{dt}
     K_\omega(h,g\,\mathrm{oneP}_\tau(X,t);Lz)
       \right|_{t=0}
     =K_\omega(h,g;L'z).
 \tag{TAE.34}
\]
Both integrals are genuinely absolutely convergent on the
original \((\nu,\mathcal F_S)\).

\emph{Proof.}
Choose a basis \(b_1,\ldots,b_d\) of \(E\), and write
\(F(h)=\sum_jF_j(h)b_j\).  For either \(L\) or \(L'\), the
kernel is the finite sum
\[
 K_\omega(h,g;L(F(h)))
          =\sum_j F_j(h)\,(\omega(h,g)Lb_j)(0).
 \tag{TAE.35}
\]
Each summand is jointly continuous by the actual
\(\mathrm{om\_cont}\) clause and continuity of \(F_j\).
Rational source invariance follows from
\(\mathrm{om\_aut}\) and the source invariance of \(F\);
rational target invariance follows from the other
\(\mathrm{om\_aut}\) clause.

The source quotient is covered by the compact set \(C_S\).
It has finite Haar volume, and every continuous source
invariant kernel is bounded by its maximum on \(C_S\).
This proves absolute convergence.  For a compact neighborhood
of a target point, (TAE.35) has a uniform bound on its product
with \(C_S\); ordinary dominated convergence proves
continuity of \(\Phi_L\).

The point derivative in (TAE.34) is TAE3 with \(p=(h,g)\)
and \(w=0\).  The one-parameter group law translates it
from zero to every \(t\).  Its derivative is therefore
\[
 K_\omega\bigl(h,g\,\mathrm{oneP}_\tau(X,t);L'(F(h))\bigr).
 \tag{TAE.36}
\]
This is jointly continuous in \((t,h)\), by (TAE.35) for
\(L'\).  On a compact interval about zero and \(C_S\) it has
a finite bound \(M\).  Rational source invariance extends
this same bound to every \(h\).
The difference quotient is dominated by \(M\), by the
one-variable integral mean-value estimate, and the constant
function \(M\) is integrable on \(\mathcal F_S\).
Dominated convergence therefore moves the derivative through
the integral.  This proves the first equality in (TAE.34).
This argument keeps the original measure and fundamental
domain; compact representatives are used only to produce a
dominating function.
% END TAE6

% BEGIN TAE7
\paragraph{TAE7: the lift is smooth and has one exponent for all real derivative words.}
\label{claim:TAE7}
The function \(\Phi_L\) satisfies the original
\(\mathrm{IsSmoothFn}(T,\theta)\) and
\(\mathrm{IsUnifModGrowth}(T,\theta)\).  In the latter
predicate one can take the same \(N\) of (TAE.2):
\[
 \forall\mathbf X,\ \forall x,\ \exists C_{\mathbf X,x}\ge0,\
 \forall g\in G_{T,\infty},\
 \bigl|R_{\mathbf X}\Phi_L(xg)\bigr|
                 \le C_{\mathbf X,x}H_T(g)^N .
 \tag{TAE.37}
\]
The derivative word consists of real Lie directions at any
of the original places, in its original order.
If a complex lowering derivative of every kernel
\(g\mapsto K_\omega(h,g;Lz)\) is zero, the corresponding
complex lowering derivative of \(\Phi_L\) is zero.

\emph{Proof.}
Repeatedly apply TAE6.  If
\(\mathbf X=((\tau_1,X_1),\ldots,(\tau_k,X_k))\) denotes
the original outer-to-inner order of \(\mathrm{foldr}\),
then
\[
 R_{\mathbf X}\Phi_L
   =\Phi_{\mathcal G_{\tau_1,X_1}\cdots
                     \mathcal G_{\tau_k,X_k}L}.
 \tag{TAE.38}
\]
Every map on the right is linear into \(\Theta\), and every
function so obtained is continuous by TAE6.

We explain why this gives the precise smoothness assertion.
On the product of the real unitary groups choose a Lie algebra
basis.  Products of its exponentials are local coordinates by
the inverse function theorem, as in TAE2.  A coordinate
derivative of a function is a finite linear combination of
its right Lie derivatives, with smooth coefficients: move
the differentiated exponential through the factors to its
right and conjugate its Lie direction by their inverse
product.  In a fixed Lie basis these adjoint coordinates
are smooth matrix functions.  TAE6 supplies the resulting
first derivatives, continuously; induction using (TAE.38)
supplies all coordinate derivatives.  Translation covers
the real group.  Composition with any finite product of
the specified one-parameter groups proves exactly the
\(\mathrm{ContDiff}\) requirement in \(\mathrm{IsSmoothFn}\),
at every adelic left translate.

For any linear \(L_1:E\to\Theta\), apply TAE5 to each
of the finitely many \(L_1b_j\), with \(C=C_S\).
Let \(M_j=\sup_{h\in C_S}|F_j(h)|\), and let \(c_j\)
be the corresponding constants in (TAE.24).
Source invariance then gives, for every \(h\),
\[
 |K_\omega(h,xg;L_1(F(h)))|
       \le\left(\sum_j M_jc_j\right)H_T(g)^N.
 \tag{TAE.39}
\]
Integration multiplies the constant by
\(\nu(\mathcal F_S)<\infty\).  Take
\(L_1=\mathcal G_{\tau_1,X_1}\cdots
\mathcal G_{\tau_k,X_k}L\).  Its constants may change
with the word, while the exponent \(N\) in TAE5 does not.
This proves the quantifier order in (TAE.37).

Finally, for a complex matrix \(Z\), the source defines
\[
 X=\tfrac12(Z+\overline Z^{\,H}),\qquad
 Y=\tfrac1{2i}(Z-\overline Z^{\,H}),\qquad
 \overline Z^{\,H}=-H_T^{-1}Z^*H_T,\qquad
 R_Z^{\mathbb C}=R_X+iR_Y .
 \tag{TAE.40}
\]
Both \(X,Y\) satisfy (TAE.13).  Apply TAE6 separately
to them.  If the sum of their pointwise kernel
derivatives is zero, its integral is zero by ordinary
linearity of the two absolutely convergent derivative
integrals.  This proves the final assertion.
% END TAE7

\subsection*{The exact lowering operator on the original Fock polynomials}

The source's real phase convention matters in the next calculation.
At a fixed place put \(\xi_\tau=\tau_E(D.\mathrm{host}.\xi_0)\).
It is nonzero and purely imaginary, with the prescribed positive
imaginary part.  On real matrix directions define
\[
 \begin{split}
 \beta_\tau(u,v)
   &=2\operatorname{Re}\!
       \left(\xi_\tau\operatorname{tr}
                 (H_S^{-1}v^*H_Tu)\right),\\
 s_\tau&=2\sqrt{2\pi|\xi_\tau|}>0.
 \end{split}
 \tag{TAE.41}
\]
This \(\beta_\tau\) is literally \(\mathrm{omDir}\) at the
same place; directions at distinct places pair to zero.
The character sign \(-1\) and the half in the original
Heisenberg law give
\[
 [D_u,D_v]=-2\pi i\,\beta_\tau(u,v)\,1 .
 \tag{TAE.42}
\]
For example, this follows by differentiating the group
commutator
\(\rho(tu)\rho(rv)\rho(-tu)\rho(-rv)
       =e^{-2\pi itr\beta_\tau(u,v)}1\).
This specifies the sign relative to the abstract coordinates
of (TAE.5); those coordinates already absorbed the original
alternating form.

In the rest of this subsection \(S\) is positive definite at
infinity.  Use its original positive basis \(B_\tau\) and the
target model's original basis \(B'_\tau=M.\mathrm{fr}_\tau^{-1}\).
They satisfy
\[
 B_\tau^*H_SB_\tau=1,\quad
 (B'_\tau)^*H_TB'_\tau
       =\operatorname{diag}(\epsilon_a),\quad
 \epsilon_a=\begin{cases}1&a<p_\tau,\\-1&a\ge p_\tau.\end{cases}
 \tag{TAE.43}
\]
The Cartan involution is
\(s=M.\mathrm{cart}_\tau
  =B'_\tau\operatorname{diag}(\epsilon_a)(B'_\tau)^{-1}\).
Thus \(J_\tau u=i s u\) is the original
\(\mathrm{cplxO}(1,s)\).  Define
\[
 A_\tau(u)=D_u+iD_{J_\tau u},\quad
 a_\tau^+(u)=s_\tau^{-1}(D_u-iD_{J_\tau u}),\quad
 u_{ac}=B'_\tau E_{ac}B_\tau^{-1}.
 \tag{TAE.44}
\]
The original polynomial vector \(\mathrm{pt}(P)\) is
\(P(a^+)f\), with one creation operator for each original
variable \((\tau,a,c)\), applied to the original vacuum \(f\).
The equality with the source's \(\mathrm{pureTensorO}\)
will be checked in TAE10.

% BEGIN TAE8
\paragraph{TAE8: a quadratic commutator determines the full polynomial action.}
\label{claim:TAE8}
Let \(a_j\) be commuting linear operators on a common invariant
complex vector space, let \(G,b_j\) be linear operators there,
and suppose
\[
 [G,a_j]=b_j,\qquad
 [b_j,a_k]=(\mu_{jk}+\mu_{kj})1,\qquad
 b_jf=0,\qquad Gf=0 .
 \tag{TAE.45}
\]
Then for every polynomial \(P\) in the finitely many variables,
\[
 GP(a)f=
    \left(\sum_{j,k}\mu_{jk}\partial_j\partial_kP\right)(a)f .
 \tag{TAE.46}
\]
No symmetry of \(\mu\), injectivity of \(P\mapsto P(a)f\), or
completion of the polynomial space is required.

\emph{Proof.}
First commute \(b_j\) across a monomial, one factor at a time.
Its final term kills \(f\), and its other terms give
\[
 b_jP(a)f=\sum_k(\mu_{jk}+\mu_{kj})(\partial_kP)(a)f .
 \tag{TAE.47}
\]
Linearity extends this from monomials to all polynomials.
Now prove (TAE.46) by induction on monomial degree.
It holds for constants because \(Gf=0\).
If it holds for \(P\), then
\[
 G(a_jP(a)f)=a_jGP(a)f+b_jP(a)f.
 \tag{TAE.48}
\]
On the polynomial side, applying
\(\sum_{k,l}\mu_{kl}\partial_k\partial_l\) to \(x_jP\)
gives \(x_j\) times its value on \(P\), plus
\(\sum_k(\mu_{jk}+\mu_{kj})\partial_kP\).
This is exactly (TAE.48), by (TAE.47).
The commuting \(a_j\)'s make every monomial evaluation
independent of its chosen order.  Addition completes the proof.
% END TAE8

% BEGIN TAE9
\paragraph{TAE9: the original vacuum is killed by the complex lowering generator.}
\label{claim:TAE9}
Fix \(\tau\).  Let \(X\) satisfy (TAE.13) and
\(sX=-Xs\), and put \(Y=isX\).  Then \(Y\) also satisfies
(TAE.13).  If \(f\in\Theta\) is an original Fock vacuum at
every real place, then
\[
       \mathcal Gf=0,\qquad
       \mathcal G=\mathcal G_{\tau,X}
                           +i\mathcal G_{\tau,Y}.
 \tag{TAE.49}
\]

\emph{Proof.}
Since \(s^*H_T=H_Ts\), the conditions on \(X\) give
\[
 Y^*H_T+H_TY
   =-iX^*H_Ts+iH_TsX
   =iH_T(Xs+sX)=0.
 \tag{TAE.50}
\]
The commutators of TAE3 imply
\[
 [\mathcal G,D_{\tau,u}]=A_\tau(Xu);
 \tag{TAE.51}
\]
the commutator is zero at every other place.
Consequently \(\mathcal Gf\) is again annihilated by
every \(A_{\tau'}(u)\): commute \(\mathcal G\) through the
two terms of \(A_{\tau'}(u)f=0\), and use
\(A_\tau(Xu)f=A_\tau(XJ_{\tau'}u)f=0\) when the places agree.
TAE3 also gives \(\mathcal Gf\in\Theta\).

Let \(k=(1,\mathrm{atTau}_\tau(s))\).
CW11 states that the canonical oscillator action of this
Cartan element is one fixed nonzero scalar on the entire
vacuum space.  Multiplication by CW24's determinant-root
character changes that scalar by another nonzero scalar,
independent of the vacuum vector.  BTH14 transfers this
statement to the actual \(\omega\).  Thus there is a single
\(c\ne0\) such that \(\omega(k)v=cv\) for every original
vacuum \(v\).  This applies both to \(f\) and to
\(\mathcal Gf\).

On the other hand,
\(s e^{tX}=e^{-tX}s\), directly from the exponential
series and \(sX=-Xs\).  TAE3 permits differentiation of
this group identity after applying \(\omega\) to \(f\):
\[
       \omega(k)\mathcal G_{\tau,X}f
                           =-c\mathcal G_{\tau,X}f .
 \tag{TAE.52}
\]
The same reasoning applies to \(Y\), since \(sY=-Ys\).
Hence \(\omega(k)\mathcal Gf=-c\mathcal Gf\).
But \(\mathcal Gf\) is a vacuum, so its left side is
\(c\mathcal Gf\).  As \(2c\ne0\), (TAE.49) follows.
This argument needs the common vacuum scalar, not a choice
of scalar depending on the tested vacuum.  CW11 and the
full-action comparison supply exactly that uniform statement.

There is also an original-data passage needed when \(f\)
is the stored \(\mathrm{IsThPVecOf}\) vector, rather than
its initial lattice vacuum \(f_0\).  We give it explicitly.
That record writes \(f\) as an \(\mathrm{IsThetaAt}\)
value of a finite split place model \(M_0\), with
vacuum \(f_0\).  Its \(\mathrm{AllSome}\) clause says that
every place model with those same data has that same
value on the prescribed finite function \(\Phi\).
The two half directions in this model have zero
infinite component.

Fix a real place and a direction \(z\).
The functional \(q\mapsto q'(0)\) is complex linear
on the subspace of functions \(\mathbb R\to\mathbb C\)
differentiable at zero.  Extend it algebraically
to a linear functional \(\ell\) on all functions.
Apply \(\ell\), pointwise, to the two real Heisenberg
curves defining \(A_\tau(z)\).  This gives a linear
operator \(\widetilde A\) on all functions on \(W\)
which agrees with \(A_\tau(z)\) on \(\Theta\).
It commutes with every purely finite Heisenberg
operator: the real and finite operators commute
before differentiation, and their fixed scalar
factors pass through \(\ell\).
Thus
\[
                  M_1=M_0+\widetilde A\,M_0
 \tag{TAE.52a}
\]
satisfies both of the same finite place-model
intertwining laws.  Its vacuum value is unchanged,
because \(\widetilde A f_0=A_\tau(z)f_0=0\).
The universal part of \(\mathrm{IsThetaAt}\) gives
\(f=M_0\Phi=M_1\Phi=f+\widetilde A f\).
Since the stored \(f\) is a theta function, this
says \(A_\tau(z)f=0\).
Apply the argument at every place and every \(z\).
It proves the literal
\(\mathrm{isFockVacAt\_thPVec}\) passage used below,
without assuming smoothness of arbitrary place-model
outputs or differentiating an arbitrary linear map.
% END TAE9

% BEGIN TAE10
\paragraph{TAE10: the actual lowering derivative is the original mixed Laplacian.}
\label{claim:TAE10}
Assume (TAE.43) at every place and let \(f\in\Theta\) be the
original vacuum at every place.  For a polynomial \(P\), suppose
at one place \(\tau\) that
\[
       \sum_{c=0}^{r_S-1}
           \partial_{\tau,a,c}\partial_{\tau,b,c}P=0
           \quad(a\ge p_\tau,\ b<p_\tau).
 \tag{TAE.53}
\]
Then for every original \(Z\in\mathfrak p^-_\tau(M)\)
and every \(h\in G_S(\mathbb A_F)\),
\[
 R_{\tau,Z}^{\mathbb C}
       \bigl(g\mapsto K_\omega(h,g;\mathrm{pt}(P))\bigr)=0 .
 \tag{TAE.54}
\]

\emph{Proof.}
The definition of \(\mathfrak p^-_\tau(M)\) says that
\(M.\mathrm{fr}_\tau ZM.\mathrm{fr}_\tau^{-1}\) is supported
only in the block with negative row and positive column.
Equivalently
\[
                         sZ=-Z,\qquad Zs=Z .
 \tag{TAE.55}
\]
Use the real parts \(X,Y\) of \(Z\) in (TAE.40).
Conjugating (TAE.55) by the Hermitian form shows
\[
        sX=-Xs,\qquad Y=isX .
 \tag{TAE.56}
\]
For example \(sH_T^{-1}=H_T^{-1}s^*\), whence
\[
 \begin{aligned}
 s(H_T^{-1}Z^*H_T)&=H_T^{-1}Z^*H_T,\\
 (H_T^{-1}Z^*H_T)s&=-H_T^{-1}Z^*H_T.
 \end{aligned}
\]
Substitution in (TAE.40) gives (TAE.56).
TAE9 now supplies \(Gf=0\) for
\(G=\mathcal G_{\tau,X}+i\mathcal G_{\tau,Y}\).

We compute the other commutators rather than identifying
them only up to a scalar.  The form in (TAE.41) satisfies
\[
 \beta_\tau(Ju,Jv)=\beta_\tau(u,v),\qquad
 \beta_\tau(u,Jv)=-\beta_\tau(Ju,v).
 \tag{TAE.57}
\]
These follow from \(s^*H_Ts=H_T\) and \(J^2=-1\).
Expanding (TAE.44) with (TAE.42) and (TAE.57) gives
\[
 \begin{split}
 [a_\tau^+(u),a_\tau^+(v)]&=0,\\
 [A_\tau(u),a_\tau^+(v)]
   &=-\frac{4\pi i}{s_\tau}
      \bigl(\beta_\tau(u,v)+i\beta_\tau(Ju,v)\bigr)\,1 .
 \end{split}
 \tag{TAE.58}
\]
At different places all these operators commute.
In particular the source's \(\mathrm{opPow}\) ordering
gives precisely \(P(a^+)f\), because the creation
operators commute.  The definition by the monomial basis
therefore agrees with \(\mathrm{pureTensorO}\) on every
polynomial.

For \(j=(\rho,a,c)\), set \(a_j=a_\rho^+(u_{ac})\).
By (TAE.51),
\[
 [G,a_j]=b_j,\qquad
 b_j=\begin{cases}
  s_\tau^{-1}\bigl(A_\tau(Xu_{ac})
                   -iA_\tau(XJu_{ac})\bigr),&\rho=\tau,\\
 0,&\rho\ne\tau.
 \end{cases}
 \tag{TAE.59}
\]
Every \(b_j\) kills \(f\).  To compute \([b_j,a_k]\),
put
\[
 o_W(a,b)=
   2\operatorname{Re}\!
       \left(\xi_\tau
              \bigl((B'_\tau)^*H_TWB'_\tau\bigr)_{ab}\right).
 \tag{TAE.60}
\]
For a real Lie element \(W\), the matrix
\((B'_\tau)^*H_TWB'_\tau\) is skew Hermitian.
Because \(\xi_\tau\) is purely imaginary,
\(o_W(a,b)=o_W(b,a)\).  If \(sW=-Ws\), its entries with
equal signs \(\epsilon_a=\epsilon_b\) vanish:
multiplication by the two diagonal sign matrices gives
\((\epsilon_a+\epsilon_b)\) times that entry equal to zero.

The original source frame produces the crucial column delta.
Since \(H_S^{-1}=B_\tau B_\tau^*\), cyclicity of trace and
the matrix units give
\[
 \beta_\tau(Wu_{ac},u_{bd})
                         =\delta_{cd}\,o_W(b,a).
 \tag{TAE.61}
\]
More explicitly the trace in (TAE.41), before taking its
real part, reduces to
\[
 \operatorname{tr}
 (E_{bd}^*(B'_\tau)^*H_TWB'_\tau E_{ac});
\]
it is zero unless \(c=d\), and otherwise is the
\((b,a)\) entry.  Thus the two original bases and the source
Hermitian inverse have all been retained.

The identities \(JXu=Yu\), \(XJu=-Yu\), and
\(JXJu=Xu\), which follow from (TAE.56), now turn
(TAE.58)--(TAE.61) into the following exact coefficients:
\[
 \begin{split}
 \nu_{ab}&=
 \begin{cases}
   2\bigl(o_X(a,b)+i\,o_Y(a,b)\bigr),
                 &a\ge p_\tau,\ b<p_\tau,\\
   0,&\text{otherwise},
 \end{cases}\\
 \mu_{(\rho,a,c),(\rho',b,d)}
   &=\begin{cases}
    -4\pi i\,s_\tau^{-2}\delta_{cd}\nu_{ab},
                         &\rho=\rho'=\tau,\\
    0,&\text{otherwise}.
   \end{cases}
 \end{split}
 \tag{TAE.62}
\]
Indeed the scalar before the factor
\(-4\pi i\,s_\tau^{-2}\delta_{cd}\) is
\[
 o_X(b,a)+i o_Y(b,a)
       -i\bigl(-o_Y(b,a)+i o_X(b,a)\bigr)
                   =2\bigl(o_X(b,a)+i o_Y(b,a)\bigr).
 \tag{TAE.63}
\]
When the two signs agree this is zero.  When they differ
it is \(\nu_{ab}+\nu_{ba}\), using symmetry of \(o_X,o_Y\).
Consequently \([b_j,a_k]=(\mu_{jk}+\mu_{kj})1\).
All hypotheses of TAE8 hold on \(\Theta\), and it proves
the exact identity
\[
 G\,\mathrm{pt}(P)=
 \mathrm{pt}\left(
   -4\pi i\,s_\tau^{-2}
       \sum_{\substack{a\ge p_\tau\\b<p_\tau}}
           \nu_{ab}\sum_c
             \partial_{\tau,a,c}\partial_{\tau,b,c}P
              \right).
 \tag{TAE.64}
\]
The inverse powers of \(s_\tau\) are legitimate by (TAE.41).
In particular (TAE.53) makes this expression zero.
Finally TAE3, followed by evaluation at zero, identifies
\(R_{\tau,Z}^{\mathbb C}K_\omega(h,\cdot;\mathrm{pt}(P))\)
with \(K_\omega(h,\cdot;G\,\mathrm{pt}(P))\).
This is zero and proves (TAE.54).
% END TAE10

% BEGIN TAE11
\paragraph{TAE11: the stored product of minors satisfies every required mixed equation.}
\label{claim:TAE11}
Let \(m=r_T,n=r_S\), and let the original integer parameters
satisfy \(p_\tau+q_\tau\le n\).
With the source's literal polynomial
\[
 \begin{split}
 \Delta_\tau={}&
 \prod_{k=1}^{p_\tau}
   \det(z_{\tau,\beta,\alpha})_{0\le\alpha,\beta<k}
                   ^{\,a^+_{\tau,k-1}-a^+_{\tau,k}}\\
 &{}\cdot
 \prod_{k=1}^{q_\tau}
   \det(z_{\tau,p_\tau+\beta,n-k+\alpha})
                      _{0\le\alpha,\beta<k}
                   ^{\,a^-_{\tau,k-1}-a^-_{\tau,k}},
 \qquad \Delta=\prod_\tau\Delta_\tau ,
 \end{split}
 \tag{TAE.65}
\]
the terminal \(a^+_{\tau,p_\tau}\) and
\(a^-_{\tau,q_\tau}\) mean zero, and the exponents are
the original natural-number differences.
Out-of-range variables have the source value zero.
For every \(a\ge p_\tau,b<p_\tau\),
\[
       \sum_c
          \partial_{\tau,a,c}\partial_{\tau,b,c}\Delta=0.
 \tag{TAE.66}
\]

\emph{Proof.}
Every variable occurring in a positive minor at \(\tau\)
has row \(<p_\tau\) and column \(<p_\tau\).
Every variable occurring in a negative minor has row
\(\ge p_\tau\) and column at least
\(n-q_\tau\ge p_\tau\).
Thus \(\Delta\) is independent of every variable whose
row and column lie on opposite sides of \(p_\tau\).
This statement survives determinants, products and all
the displayed nonnegative powers.  Variables at other
places are distinct.

For a given column \(c<p_\tau\), the derivative
\(\partial_{\tau,a,c}\Delta\) is zero.
For \(c\ge p_\tau\), the derivative
\(\partial_{\tau,b,c}\Delta\) is zero.
Polynomial partial derivatives commute, as is verified
on monomials, so each summand of (TAE.66) is zero.
Summing gives the result.  No monotonicity or irreducibility
assertion about the polynomial representation is needed
for this harmonicity statement.

For the original bundle \(p_\tau+q_\tau=r_T\le r_S\).
Consequently TAE10 applies to the actual
\(\mathrm{RealDat.Delta}(i)\), with its original
Fock normalization, at every parameter \(t\).
% END TAE11

\subsection*{The exact interface to lowest-type and Hilbert realization}

% BEGIN TAE12
\paragraph{TAE12: lowering propagates to all coefficient values under the stated lowest-type hypotheses.}
\label{claim:TAE12}
Let \(c:E\) be the original finite-dimensional coefficient
representation with \(c.\mathrm{Is}\).  Assume explicitly
\(\mathrm{LowSm}(M)\) and \(\mathrm{LowIrr}(M)\), and let
\[
 io:\mathrm{LowT}(M)\longrightarrow
              \operatorname{Hom}_{\mathbb C}(E,\mathcal P)
 \tag{TAE.67}
\]
satisfy the original \(\mathrm{IsIoW}\) with this
\(\omega,\mathrm{pt},\Delta\).
If all kernel lowering derivatives of
\(\mathrm{pt}(\Delta)\) vanish, then for every pure
lowest-type tensor \(v\), every \(w\in E\), and every source
point \(h\), all lowering derivatives of
\[
             g\longmapsto
             K_\omega(h,g;\mathrm{pt}(io(v)w))
 \tag{TAE.68}
\]
vanish.  Its integral has the same property.

\emph{Proof.}
For precision, \(\mathrm{IsIoW}\) supplies three things.
First, for every real source element \(u\),
\[
       \omega(u,1)\mathrm{pt}(io(v)w)
                           =\mathrm{pt}(io(v)c(u)w).
 \tag{TAE.69}
\]
Second, a compact target element acting in one model
factor replaces that factor of \(v\), with its literal
frame-conjugated adelic lift on the \(\omega\) side.
Third, there are unit vectors \(v_{0,\tau}\) and
\(w_0\) with \(io(\otimes_\tau v_{0,\tau})w_0=\Delta\).
We use all three clauses.

Let \(\mathcal N\subset\mathcal P\) consist of polynomials
whose kernels are killed by every lowering derivative,
at every source point.  It is a complex vector subspace.
Indeed TAE3 and the two-real-part formula (TAE.40) express
any lowering derivative as a fixed linear operator on
the theta vector followed by \(\omega(h,g)\) and
evaluation.  Thus addition and scalar multiplication
preserve the vanishing.

For the fixed tensor \(v_0=\otimes_\tau v_{0,\tau}\),
the subspace
\(\{w:io(v_0)w\in\mathcal N\}\) contains \(w_0\ne0\).
It is invariant under every \(c(u)\).  For a real source
element, (TAE.69) replaces the kernel source variable
\(h\) by \(hu\), and this does not affect a condition
imposed for all \(h\).  A finite source element acts
trivially by \(c.\mathrm{Is}\).
Every adelic element is a product of its real-supported
and finite-supported parts; the original infinity chart
identifies its real part with the same real group.
The irreducibility clause in \(c.\mathrm{Is}\) therefore
makes this subspace all of \(E\).

Now fix \(w\).  The inverse image of \(\mathcal N\)
under \(v\mapsto io(v)w\) contains \(v_0\) and is stable
under each compact target factor.
Here is the derivative check in that assertion.
For a real target point \(k\),
\[
 R_Z^{\mathbb C}(r_k\varphi)
        =r_k R_{\operatorname{Ad}(k^{-1})Z}^{\mathbb C}\varphi,
       \qquad (r_k\varphi)(g)=\varphi(gk).
 \tag{TAE.70}
\]
It follows by differentiating
\(g e^{tX}k=gk e^{t k^{-1}Xk}\), separately in the
two real parts.  Smoothness needed here was proved
by TAE3 and TAE7.  In the model frame a compact
element commutes with \(\operatorname{diag}(\epsilon_a)\),
so it is block diagonal.  Conjugation by it preserves
the negative-row, positive-column block
\(\mathfrak p^-_\tau(M)\).  Elements at other places
leave that block unchanged.  The target clause of
\(\mathrm{IsIoW}\) and (TAE.70) prove the claimed stability.

By \(\mathrm{LowSm}\), each full lowest compact type lies
inside the model's smooth space.  By \(\mathrm{LowIrr}\),
it is a finite-dimensional irreducible compact
representation.  The algebraic span of the orbit of
each nonzero \(v_{0,\tau}\) is therefore its whole
lowest type.  Apply this fact in one tensor factor
at a time.  Multilinearity proves that the entire
algebraic tensor product maps into \(\mathcal N\).
This gives (TAE.68) for every stated \(v,w,h\).
Finally apply the last assertion of TAE7 to pass
the vanishing through the original source integral.

The two assumptions \(\mathrm{LowSm}\) and
\(\mathrm{LowIrr}\) have been used explicitly.
\(\mathrm{DSModel.ds}\) alone supplies a lowest
irreducible compact subspace; it does not definitionally
say that the supremum of all such subspaces is one copy.
This proof does not claim to have supplied that
multiplicity-one argument.
% END TAE12

% BEGIN TAE13
\paragraph{TAE13: the original lowest-type integral now has the complete analytic and cusp packet, with its remaining producers visible.}
\label{claim:TAE13}
Fix one original \(\mathrm{RealDat}\ R\) on the same bundle.
Retain the literal hypotheses
\(R.\mathrm{toLevCh.Is}(A.\mu_0)\),
\(R.\mathrm{IsModels}\),
\(R.\mathrm{thP\_\omega P}\),
\(R.\mathrm{thP\_vP}\), and
\(R.\mathrm{thP\_pt\_theta}\).
These keep its original full \(\omega_P(i,t)\),
vacuum \(v_P(i,t)\), polynomial map and source models.
For a fixed \(i\), additionally
suppose that one map \(io_i\), independent of \(t\), satisfies
\[
 \forall t,\quad
 \mathrm{IsIoW}\bigl(
     R.\theta,\ R.MW(i,\mathrm{sg}_0),\ R.cV(i),\
     R.\omega_P(i,t),\ R.\mathrm{ptP}(i,t),\
     \mathrm{Delta}(i),\ io_i\bigr),
 \tag{TAE.71}
\]
and retain \(\mathrm{LowSm}\) and \(\mathrm{LowIrr}\)
for that same target model.
Put \(F_{i,t,x}=R.\mathrm{extV}(i,t)x\) and
\(L_v=R.\mathrm{ptP}(i,t)\circ io_i(\otimes_\tau v_\tau)\).
Then, for every \(t\) and original source Hilbert vector \(x\),
its literal lowest-type integral
\[
 U_{i,t,x}(v)(g)=
  \int_{R.\mathrm{fdV}(i,t)}
   K_{R.\omega_P(i,t)}\bigl(h,g;L_v(F_{i,t,x}(h))\bigr)
                       \,dR.\nu_V(i,t)(h)
 \tag{TAE.72}
\]
satisfies the exact \(\mathrm{IsHolCuspKMap}\) analytic
packet at the original principal target level.
It is linear in \(x\) and multilinear in the lowest-type
vectors.  Its uniform growth exponent can be chosen as
\[
       N_i=2\,r_{V_c(i)}r_{W_c(i)}\,\#\mathrm{Tau},
 \tag{TAE.73}
\]
independently of \(t,x,v\) and of the real derivative word.
This assertion does not construct an element of
\(\mathrm{MMap}(R.MW(i,\mathrm{sg}_0))\).

\emph{Proof.}
The source in (TAE.72) is the same finite-dimensional
coefficient-valued extension used in BTH3.  Its continuity
and rational invariance are the original coefficient-model
identities, and its source quotient has BTH3's compact
representatives.  The map \(L\) in (TAE.4) is, for fixed
\(v\), the literal
\[
             L=R.\mathrm{ptP}(i,t)\circ
                          io_i(\otimes_\tau v_\tau).
 \tag{TAE.74}
\]
TAE6 proves absolute convergence, continuity and rational
target invariance.  TAE7 proves smoothness and uniform
moderate growth with (TAE.73), since the real Hom space
has exactly that dimension.  This exponent is selected
before the coefficient map, the vector and the derivative
word; only their bounding constants change.

The original finite covariance and principal-level fixing
were proved in BTH5 and BTH15, with the same action and
the same frames.  Only the finite vacuum calculation in
those proofs is used here; its subsequent extension to a
full multiplicity map is not used.  For a principal
target element \(k\), take its matched source element
to be \(1\).  Every reduced block is the identity, so
all determinant and half-support multipliers in
BTH.19 are one.  Thus \(\omega(1,k)v_P=v_P\).
Finite action commutes with every real creation
operator by the \(L^2\) derivative transport of TAE5,
so it fixes every original polynomial vector.
This calculation is pointwise in \(t\), uses only
the stated level and vacuum fields, and applies
unchanged at each \(t\).  The original principal
target character is one.  Inserting this identity under
the integral gives the required right level law.
The compact target clause of (TAE.71), inserted under
the same integral, gives its exact model-frame compact
equivariance.  Multilinearity and linearity in \(x\)
follow from linearity of the kernel and of the original
coefficient extension, using the absolute convergence
already established.

The original rank inequality \(r_W\le r_V\) and signature
order give the hypotheses of TAE11 for
\(\mathrm{Delta}(i)\).  The actual \(\mathrm{IsThPVecOf}\)
and polynomial-theta fields give the Fock-vacuum equations
for \(v_P(i,t)\), by the explicit
\(\mathrm{AllSome}\) argument at the end of TAE9.
It changes neither its scale nor its polynomial map.
TAE10 and TAE11 kill the lowering kernel of
\(\mathrm{ptP}(i,t)(\mathrm{Delta}(i))\).
TAE12 propagates this to every value of (TAE.74), and
TAE7 passes it through the integral.

FIA8 gives the finite isotropic average for every
polynomial and every source and target translate.
CTC9--CTC13 therefore give the zero transvection
constant terms and genuine full cuspidality of
(TAE.72).  Their unipotent quotient integrals converge;
no undefined-integral convention is used here.
Together these are exactly the eight fields
of the source's \(\mathrm{IsHolCuspKMap}\):
continuity, rational left invariance, level law,
smoothness, uniform growth, cuspidality, compact
equivariance and lowering.

Finally (TAE.71) is the original \(\mathrm{PIoP}(i,io_i)\).
It is not replaced by the weaker quantifier order
\(\forall t\,\exists io_{i,t}\).
If the original \(R.\mathrm{ioP}(i)\) is its stored
\(\mathrm{OrZero}\) choice, existence of a witness
with (TAE.71) activates that choice by the definition
of \(\mathrm{OrZero}\).  One may then use the actual
stored map in (TAE.72).
Nothing in this implication supplies such a witness,
or supplies \(\mathrm{LowSm}/\mathrm{LowIrr}\).
Nothing in it extends \(U_{i,t,x}\) from the lowest
compact type to all smooth vectors of the target
discrete-series model.  Those are the remaining
interfaces, stated next.
% END TAE13

\subsection*{The exact remaining producer}

The analytic part just reconstructed is pointwise on one
supplied original \(R\).  It introduces no new family
threshold, coefficient field, Hermitian adjoint, or theta
matrix.  At the base tuple it uses the same
\(t_0=A_3.A.\mu_0\) as BTH12 and SLR8.  Its consequences
can also be applied at all \(t\), with the common
polynomial map quantified as in (TAE.71).

The residual common-object obligations are as follows.
\begin{enumerate}
\item Produce one \(io_i\) satisfying (TAE.71), with both
original intertwining laws and the unit-vector
normalization to the literal \(\mathrm{Delta}(i)\).
The source supplier is
\texttt{JointHarm.exists\_\allowbreak io\_\allowbreak poly}, followed by
\texttt{JointHarm.exists\_\allowbreak ioP}.
Its finite polynomial lowest-weight transfer has not
been independently reconstructed in this chapter.
\item Prove \(\mathrm{LowSm}\) and
\(\mathrm{LowIrr}\) for the full original
\(\mathrm{lowKSub}\).
The source supplier is
\begin{quote}\small
\texttt{HolLowOnce.lowSm\_and\_lowIrr}.
\end{quote}
In particular its multiplicity-one step uses a
holomorphic coefficient uniqueness argument.
The earlier nonzero smooth-vector and cyclicity
argument in BTH1 is not a substitute for that step.
\item Supply boundedness for the actual smooth cuspidal
uniform-moderate-growth lowest-type form, and the
finite-volume target quotient needed by the Hilbert
realization.  The source route uses
\(\mathrm{Published.CuspRapidDecay}\) through
\texttt{bounded\_\allowbreak of\_\allowbreak smooth\_\allowbreak
cuspidal\_\allowbreak uniform\_\allowbreak moderate}.
Naming this predicate does not certify a published
application: its hypotheses must still be matched,
or the needed instance proved directly.
\item Reconstruct the actual bounded holomorphic
coefficient identity and smooth multiplicity-map
realization.  The immediate suppliers are
\begin{quote}\small
\texttt{inner\_reg\_eq\_coef\_extremal},\\
\texttt{boundedHolomorphicForm\_eq\_multiplicitySpace\_value}.
\end{quote}
The remaining conclusion is an actual
\[
 \Theta_{i,t,x}\in\mathrm{MMap}(R.MW(i,\mathrm{sg}_0))
\]
satisfying \(\mathrm{IsMultMap}\), whose lowest-type
restriction is exactly (TAE.72), with its original
Sobolev continuity, finite level and central-finiteness.
\end{enumerate}
These are retained C producer subcontracts in the present
audit, not four new independent top-level obligations.
The proved implications TAE1--TAE13 are B.
No internal D translation issue is being concealed by
calling any of these producers automatic.
The geometric Hilbert reading, the simultaneous rational
tests and the original base-row producer in BTH12/SLR8
remain separate from this analytic extension.
